\chapter{Introduction}
\label{chap:intro}

Swarm robotics studies how many relatively simple agents can coordinate through local interaction without continuous centralized control~\cite{diasSwarmRoboticsPerspective2021,debieSwarmRoboticsSurvey2023}. In practice, that coordination depends on message exchange over dynamic wireless links, so communication reliability and communication security become intertwined design concerns rather than separate subsystems~\cite{varadharajanSwarmRelaysDistributed2020,chenSecuringEmergentBehaviour2021}.

The literature reviewed in this thesis suggests a recurring trade-off. Lightweight cryptographic mechanisms are attractive for constrained UAV and IoT-style platforms because they reduce key size, messaging overhead, and runtime cost~\cite{adeniyiSystematicReviewECC2024,hanLightweightSecureCommunication2025,manikandanASMTPAnonymousSecureMessaging2024}. At the same time, stronger decentralized trust mechanisms, such as blockchain-assisted coordination or attestation frameworks, can improve system assurance but usually add latency, bandwidth cost, or implementation complexity~\cite{sunResearchDataSecurityEmergency2022,pachecoSecuringFederatedLearning2024,kohliSwarmNetFirmwareAttestation2024}. Within this representative corpus, adaptive control of cryptographic settings is discussed less often than the underlying primitives themselves.

This thesis addresses that narrower gap by developing and evaluating an adaptive lightweight communication framework that dynamically balances security-oriented settings and performance-oriented settings in a swarm-inspired environment. The relevance of the topic follows from the growing use of distributed robotic and IoT-style systems in settings where communication must remain protected even when energy, bandwidth, and processing capacity are limited.

The purpose of the thesis is to design and evaluate a reproducible adaptive secure-communication framework for swarm-oriented systems. The main tasks are to review relevant swarm-security and lightweight-cryptography literature, define a formal model and threat model, specify adaptive cryptographic profiles, implement a pairwise benchmark prototype, and compare static-heavy, static-balanced, static-lightweight, and adaptive operation under constrained execution.

The object of study is secure communication in swarm-oriented distributed systems. The subject of study is adaptive selection of cryptographic profiles for the protected data plane. The scientific novelty of the work is the combination of a bounded formal model with an executable benchmark that exposes the security--performance trade-off instead of treating the adaptive framework only as a conceptual architecture.

The swarm network is modeled as a dynamic graph $G(V,E,t)$, and the adaptation function
\begin{equation}
F : S(t) \rightarrow P,
\label{eq:adaptation_function}
\end{equation}
maps swarm-state variables $S(t)$ to configurable cryptographic parameters $P$. At the conceptual level, $S(t)$ includes density, energy, packet loss, and threat estimates, while $P$ includes key length, cipher selection, authentication mode, and key-rotation frequency. At the executable level, the benchmark prototype instantiates a smaller pairwise subset of this idea: it switches between three runtime profiles using scripted threat and energy inputs and evaluates the resulting data-plane behavior under constrained CPU and memory budgets.

The contribution is therefore twofold. First, the thesis develops a formal and architectural model for adaptive secure communication in swarm-oriented systems. Second, it grounds that model in a reproducible pairwise benchmark prototype instead of relying only on analytical estimates. The implementation artifacts, scenario presets, tests, and final benchmark matrix are organized as a reproducible repository artifact accompanying the thesis. The resulting evidence is intentionally narrower than a field deployment: it supports claims about the behavior of a constrained host-based prototype, not about full multi-hop swarm operation on physical robots.
